% Independent analytical audit; input fragment.
\section{A finite graph hierarchy for the original reference potential wave}
\label{sec:original-gauge-graph}

This section concerns the original-reference pure-gauge normal model, not the
conformal-harmonic counterexample operator.  It provides an a priori graph
estimate and a strong-normal embedding for its homogeneous scalar sector.
Existence by closure of a common invariant graph domain, the full $S\ne0$
elliptic boundary coupling, and transfer to variable CMC coefficients remain
separate tasks.

\subsection{Exact forced scalar equation}

In the null-halfspace notation of Section~\ref{sec:gauge-counterexamples},
let $S_i$ be the spatial covariant York generator, satisfying
$\Delta S_i+\partial_i\operatorname{div}S/3=0$, and write
$\xi^t=S_x-x\psi$, $\xi^x=x\psi$.  Direct physical diffeomorphism gives
\begin{align*}
 h&=-L_{\rm conf}S,&
 c&=\tfrac23\operatorname{div}S-2\psi,\\
 u&=-D S_x+xD\psi,&
 v&=-DS+\nabla(S_x-x\psi),\\
 k&=x\Delta S_x-3\partial_xS_x-x^2\Delta\psi+x\psi_x+3\psi.
\end{align*}
For the original reference-lapse pole $(-q-v_x-u)/x$, with $q=k$ in
the constraint-free sector, substitution gives the exact equation
\[
 P\psi=\frac{\mathcal L S_x}{x}+\frac{2\partial_tS_x}{x^2},\qquad
 \mathcal L=\partial_t^2-2\partial_t\partial_x-\Delta_\perp,
 \qquad P=\mathcal L+\frac2{x^2}.
\]
The singular source has an explicit regular lift.  Since
$[\mathcal L,x]=-2\partial_t$, the substitution
\[
 \psi=\partial_t S_x+\frac{x}{2}\mathcal L S_x+\rho
\]
gives
\begin{equation}
 P\rho=-\frac{x}{2}\mathcal L^2 S_x.
 \label{eq:original-reference-forced-lift}
\end{equation}
Thus smooth York data produce an $O(x)$ bulk source after the two displayed
boundary jets are lifted.  Estimating the York generator through its coupled
boundary evolution is still required; this identity does not assume that
source is independently prescribed.

\subsection{Null energy and spatial dilation}

First consider $P\rho=0$.  On a finite normal collar, or with suitable decay at
the inner end, define
\[
 E[\phi]=\int\left(|\phi_t|^2+|\nabla_\perp\phi|^2+
 \frac{2|\phi|^2}{x^2}\right)\,dx\,dy\,dz.
\]
For smooth regular fields with justified traces, multiplication of
$P\phi=F$ by $2\phi_t$ gives
\[
 \frac{d}{dt}E[\phi]
 =-2\int_{x=0}|\phi_t|^2\,dy\,dz
 +2\operatorname{Re}\int\overline{\phi_t}F,
\]
up to the separately controlled inner-boundary flux.  Consequently
$d\sqrt E/dt\leq\|F\|_{L^2}$.  The energy is degenerate in the normal
derivative; positive potential does not itself supply $\phi_x$.

Use the spatial-time dilation
\[
 B=t\partial_t+x\partial_x.
\]
Its exact commutator is
\[
 [P,B]=2P+2\Delta_\perp.
\]
In particular $P(B\rho)=2\Delta_\perp\rho$.  More generally,
\[
 P B^j\rho=\Delta_\perp\big[(B+2)^j-B^j\big]\rho.
\]
Time and tangential derivatives commute with $P$.  The finite hierarchy
\begin{equation}
 \mathcal G_N(t)=
 \sum_{j+p+|\alpha|\leq N}
 E[B^j\partial_t^p\partial_\perp^\alpha\rho](t)
 \label{eq:original-reference-graph-norm}
\end{equation}
therefore closes: every source contains at most $j-1$ dilation operators and
two tangential derivatives, and its $L^2$ norm is controlled by an energy at
the same index using that energy's tangential derivative.  A finite sum and
Cauchy--Schwarz give
\[
 \mathcal G_N(t)\leq C_{N,T}\mathcal G_N(0),\qquad 0\leq t\leq T,
\]
for this a priori regular class.  The same argument permits independently
bounded commuted sources.  It does not hide a normal derivative in the source.

\subsection{Strong-normal embedding}

The elementary identity
\[
 \phi_x=\frac{B\phi-t\phi_t}{x}
\]
implies
\[
 \|\phi_x\|_{L^2}\leq\frac1{\sqrt2}
 \left(\sqrt{E[B\phi]}+|t|\sqrt{E[\phi_t]}\right).
\]
Thus the finite hierarchy controls a full, unweighted normal derivative, not
merely a conormal derivative.

For any commuted field $P\phi=F$, the equation also gives
\[
 \frac{2\phi}{x^2}=F-\phi_{tt}+2\phi_{tx}+\Delta_\perp\phi.
\]
Apply the previous normal estimate to $\phi_t$.  This controls
$\|\phi/x^2\|_{L^2}$ by finitely many energies with two additional
time/dilation/tangential indices and the known source.  No assertion that a
Hardy inequality alone controls the inverse-square potential is needed.

For example,
\[
 \frac{B^2\rho}{x^2}=
 \rho_{xx}+\frac{\rho_x}{x}
 +\frac{2t\rho_{tx}}x
 +\frac{t^2\rho_{tt}}{x^2}+\frac{t\rho_t}{x^2},
\]
while
\[
 \frac{\rho_x}{x}=\frac{B\rho-t\rho_t}{x^2},\qquad
 \frac{\rho_{tx}}x=\frac{B\rho_t-t\rho_{tt}}{x^2}.
\]
The inverse-square estimates applied to these finitely many commuted fields
give an unweighted $\rho_{xx}$ bound from $\mathcal G_4$ on finite time
intervals.  Similarly $B^3\rho/x^2$ controls $x\rho_{xxx}$ once the lower
terms are estimated.  Thus $\mathcal G_5$ controls the basic pure-gauge
geometric fields
\[
 c=-2\rho,\quad u=xD\rho,\quad v=-\nabla(x\rho),\quad
 a=-[x\nabla\nabla\rho+G\otimes\nabla\rho+
 \nabla\rho\otimes G]^{\rm TF},
\]
with $c$ in $H^2$, $u,v,a$ in $H^1$, and
$k=-x^2\Delta\rho+x\rho_x+3\rho$ in $H^1$, on a bounded collar.
These are finite graph estimates with explicitly stronger initial graph
regularity; they are not estimates in an unchanged minimal NR Sobolev norm.
The corresponding estimates for a general nonzero York lift require its
commuted source in~\eqref{eq:original-reference-forced-lift} to be controlled.

\subsection{Why powers of the time generator alone do not suffice}

At a fixed interior point, suppress tangential frequencies and freeze the
positive potential at $V_0$.  The first-order time generator on $(\rho,u)$
has normal symbol
\[
 \begin{pmatrix}0&1\\-V_0&2i\xi\end{pmatrix},\qquad
 s_\pm=i\xi\pm i\sqrt{\xi^2+V_0}.
\]
For positive large $\xi$, its slow branch is
$s_-=V_0/(2i\xi)+O(\xi^{-3})$, an operator of order $-1$ rather than
a normally elliptic operator.  Thus a claim that finitely many powers of the
time generator automatically control normal Sobolev derivatives is not
justified.  Microlocal slow-branch packets expose this obstruction.  The
independent dilation commutations in~\eqref{eq:original-reference-graph-norm}
are the constructive replacement.

The estimates above are for smooth regular solutions, or limits for which
the flux and commutator calculations are established.  Proving existence on
a common invariant completed graph domain, handling its trace conditions,
and closing the coupled elliptic York boundary dynamics are necessary before
promoting this scalar a priori result to a full well-posedness theorem.
